%% ==============================================================
%% КОНФИГУРАЦИЯ — Вариант 19
%% Геометрия B, тип 9 (Комбинированный двухуровневый)
%% Сгенерировано автоматически — не редактировать вручную
%% ==============================================================

% --- Информация о студенте ---
\newcommand{\varNumber}{19}
\newcommand{\varStudent}{Строгов~М.М.}
\newcommand{\varGroup}{ЗНБ21-02Б}
\newcommand{\varSupervisor}{Башмур~К.А.}

\newcommand{\varStudentID}{82159416}

% --- Геометрия подшипника (геометрия B) ---
\newcommand{\varR}{0.05}
\newcommand{\varRmm}{50}
\newcommand{\varC}{7e-05}
\newcommand{\varCmkm}{70}
\newcommand{\varL}{0.08}
\newcommand{\varLmm}{80}
\newcommand{\varLD}{0.8}
\newcommand{\varGeomLetter}{B}

% --- Тип микрорельефа (тип 9) ---
\newcommand{\varDimpleType}{9}
\newcommand{\varDimpleTypeName}{Комбинированный двухуровневый}
\newcommand{\varDimpleTypeGenitive}{комбинированным}

% --- Параметры микрорельефа ---
\newcommand{\varDimpleA}{0.00248}
\newcommand{\varDimpleAmm}{2.48}
\newcommand{\varDimpleB}{0.00228}
\newcommand{\varDimpleBmm}{2.28}
\newcommand{\varDimpleHp}{1e-05}
\newcommand{\varDimpleHpmkm}{10}
\newcommand{\varDimpleHpRatio}{0.1}

% --- Формула зазора ---
\newcommand{\varDimpleFormula}{
    \Delta h = h_p \text{ при } r\leq0.5;\; 0.6h_p \text{ при } 0.5<r\leq0.8
}

% --- Общие параметры ---
\newcommand{\varN}{2980}
\newcommand{\varEta}{0.01105}
\newcommand{\varNphi}{8}
\newcommand{\varNZ}{11}
\newcommand{\varPhiRange}{$90^\circ \div 270^\circ$}
